% ======================================================================
% Kapitel 2 — Background and Related Work (Budget: ~2,5 Seiten, KURZ)
% Struktur nach Salmas Kapitel 2 (Erkenntnisse + Luecken), plus eine
% kompakte Foundations-Section, weil CTI-/RAG-Begriffe nicht
% selbsterklaerend sind. Keine Designkriterien, keine Evaluations-
% methodik hier --- das lebt in Kap. 3 (\ref{sec:requirements}, \ref{sec:evaluation-design}).
% ======================================================================
\chapter{Background and Related Work}
\label{ch:background}

\section{CTI and RAG foundations}
\label{sec:cti}
% B1--B3 speisen die Gaps (\ref{sec:research-gap}) und Kriterien
% (\ref{sec:requirements}); Rueckverweis in \ref{sec:assessment}.

Cyber Threat Intelligence (CTI) is evidence-based knowledge about existing
or emerging threats, including context, indicators, and actionable advice
that informs security decisions \parencite{sun2023ctimining}. In practice,
CTI is communicated through recurring artifact types. Indicators of
Compromise (IOCs) are forensic artifacts, such as file hashes, IP
addresses, or domains, that indicate a compromised system
\parencite{palacin2021threat}. Tactics, Techniques, and Procedures (TTPs)
describe adversary behavior on three levels: the tactical goal, the method
used to achieve it, and its concrete implementation
\parencite{ismail2025wazuh}. The MITRE ATT\&CK framework is the industry
standard for categorizing these behaviors \parencite{lekssays2025techniquerag}.
Threat hunting builds on CTI: analysts proactively and iteratively search
their environment for signs of compromise that automated detection has
missed, typically starting from intelligence about the adversary
\parencite{palacin2021threat,sun2023ctimining}.

Analyst workflows follow a recurring cycle from planning and collection
through processing and analysis to dissemination
\parencite{palacin2021threat}. Across this cycle, the literature reports
three bottlenecks:
% TODO: sobald die Workflow-Abbildung vorliegt, hier einen Verweissatz
% auf Figure~\ref{fig:analyst-workflow} ergaenzen.

% TODO Abbildung: generischer CTI-Analyst-Workflow (Feed-/Report-Eingang
% -> Triage -> Anreicherung/Korrelation -> Bewertung/Report), mit
% markierten Engpaessen B1--B3. Als drawio anlegen (docs/architecture/).
% \begin{figure}[htbp]
%     \centering
%     \includegraphics[width=0.9\textwidth]{cti_analyst_workflow}
%     \caption{Generic CTI analyst workflow with bottlenecks B1--B3}
%     \label{fig:analyst-workflow}
% \end{figure}

\begin{enumerate}
    \item \textbf{B1 --- Information overload:} the volume, heterogeneity,
          and informal structure of threat feeds and reports exceed what
          analysts can process consistently
          \parencite{alhuzali2025llm,rastogi2023soc}.
    \item \textbf{B2 --- Recency:} threat knowledge changes faster than
          static knowledge bases, including the parametric knowledge of
          trained language models, can follow
          \parencite{paul2025proactive,chen2023benchmarking}.
    \item \textbf{B3 --- Verification effort:} extracting information from
          unstructured reports and validating it remains manual expert
          work, and outputs without transparent provenance cannot be acted
          on \parencite{rajapaksha2024rag,tseng2024llm}.
\end{enumerate}

Retrieval-Augmented Generation (RAG) grounds LLM output in external
knowledge. A RAG pipeline indexes documents by cleaning, chunking, and
embedding them, retrieves the passages most relevant to a query, optionally
reranks them, and assembles them into the prompt from which the model
generates its answer \parencite{gao2024rag}. This design targets the LLM
limitations most relevant to CTI: hallucination, outdated knowledge, and
reasoning that cannot be traced to sources
\parencite{gao2024rag,chen2023benchmarking}.
% TODO: sobald die Pipeline-Abbildung vorliegt, Verweissatz auf
% Figure~\ref{fig:rag-pipeline} ergaenzen.

Retrieval can be sparse, dense, or hybrid;
Table~\ref{tab:retrieval-strategies} contrasts the strategies and the role
of reranking \parencite{agnihotram2025retrieval}.

\begin{table}[H]
    \centering
    \small
    \begin{tabular}{lp{0.26\textwidth}p{0.24\textwidth}p{0.24\textwidth}}
        \toprule
        \textbf{Strategy} & \textbf{Principle} & \textbf{Strength} & \textbf{Limitation} \\
        \midrule
        Sparse (BM25) & exact term matching with term-frequency weighting
            & efficient; reliable for explicit identifiers
            & misses semantically equivalent formulations \\
        Dense & embeddings in a shared vector space
            & captures paraphrases and related concepts
            & can surface related yet irrelevant passages \\
        Hybrid & parallel sparse and dense retrieval, merged by rank or
                 score fusion
            & keeps exact matches while adding semantic coverage
            & merged list still needs reordering \\
        Reranking & cross-encoder scores each query--passage pair
            & raises precision with little recall loss
            & adds latency \\
        \bottomrule
    \end{tabular}
    \caption{Retrieval strategies and reranking in RAG pipelines}
    \label{tab:retrieval-strategies}
\end{table}

These trade-offs matter for CTI text, where the same threat actor appears
under aliases such as APT28 and Fancy Bear and distinct ATT\&CK techniques
share overlapping keywords; hybrid retrieval with reranking therefore
suits CTI corpora \parencite{alhuzali2025llm,lekssays2025techniquerag}.
Citing the retrieved sources in the generated answer makes it verifiable:
analysts can inspect the evidence behind each claim instead of trusting
the model \parencite{rajapaksha2024rag,liu2023trustworthy}.

% TODO Abbildung: generische RAG-Pipeline (Ingestion -> Chunking ->
% Indexing -> Retrieval -> Reranking -> Context Assembly -> Generation).
% Ersetzt Prosa-Definitionen; 3.2 zeigt spaeter die konkrete Instanz.
% \begin{figure}[htbp]
%     \centering
%     \includegraphics[width=\textwidth]{rag_pipeline_generic}
%     \caption{Generic RAG pipeline stages}
%     \label{fig:rag-pipeline}
% \end{figure}

\section{Key findings from the literature}
\label{sec:literature-findings}

The literature review covers four topic areas: CTI analyst workflows, LLMs
in cybersecurity, RAG for CTI, and the evaluation of RAG systems.
Table~\ref{tab:literature-findings} summarizes the key findings; the
research gaps in Section~\ref{sec:research-gap} follow from them.

\begingroup
\small
\begin{longtable}{p{0.15\textwidth}p{0.53\textwidth}p{0.22\textwidth}}
        \toprule
        \textbf{Topic area} & \textbf{Key finding} & \textbf{Source} \\
        \midrule
        \endfirsthead
        \toprule
        \textbf{Topic area} & \textbf{Key finding} & \textbf{Source} \\
        \midrule
        \endhead
        \midrule
        \multicolumn{3}{r@{}}{\footnotesize\itshape continued on next page} \\
        \endfoot
        \bottomrule
        \caption{Key findings of the literature review}
        \label{tab:literature-findings}
        \endlastfoot
        CTI analyst workflows
        & The volume and the non-uniform, informal structure of threat
          feeds overwhelm analysts; fully automated processing is not
          feasible & \cite{rastogi2023soc,alhuzali2025llm} \\
        & Analyzing natural-language CTI reports is labor-intensive and
          repetitive and delays response
          & \cite{tseng2024llm,rajapaksha2024rag} \\
        \midrule
        LLMs in cybersecurity
        & LLMs support vulnerability analysis, report summarization, and
          IOC/TTP extraction; junior analysts gain the most
          & \cite{divakaran2024llm} \\
        & Recurring failure modes are hallucination, outdated knowledge,
          and technical misinterpretation & \cite{alhuzali2025llm} \\
        & LLMs over-trust retrieved context and can be misled by incorrect
          information in it & \cite{chen2023benchmarking} \\
        \midrule
        RAG for CTI
        & Hybrid retrieval with reranking addresses terminology variance
          (actor aliases) and overlapping keywords in security text
          & \cite{alhuzali2025llm,lekssays2025techniquerag} \\
        & Retrieval grounding raises precision in TTP mapping, where
          decoder-only LLMs show high recall but low precision
          & \cite{fayyazi2024ttp} \\
        & Existing CTI-RAG prototypes each target a single task, e.g.\
          attack investigation, technique annotation, event response, or
          IOC extraction
          & \cite{alhuzali2025llm,lekssays2025techniquerag,ismail2025wazuh,tseng2024llm} \\
        \midrule
        RAG evaluation
        & Retrieval is assessed with Precision@k and Recall@k plus
          rank-aware metrics such as MRR and nDCG
          & \cite{agnihotram2025retrieval,yu2024evaluation} \\
        & Reference-free frameworks (RAGAS, ARES) score faithfulness,
          answer relevancy, and context relevance without gold answers
          & \cite{es2023ragas,saadfalcon2024ares} \\
        & LLM-as-a-judge scoring risks reinforcing model biases and can
          diverge from human judgment in specialized domains
          & \cite{akkiraju2024facts,yu2024evaluation} \\
        & Industrial ground-truth datasets for CTI tasks are scarce; only
          few public benchmarks (e.g.\ CTIBench) exist
          & \cite{xu2025llm,alhuzali2025llm} \\
\end{longtable}
\endgroup

\section{Research gaps}
\label{sec:research-gap}

From the findings in Table~\ref{tab:literature-findings}, three research
gaps are derived:

\begin{enumerate}
    \item \textbf{G1 --- Workflow-level empirical evidence:} Existing
          CTI-RAG prototypes each address a single task, such as attack
          investigation or technique annotation. How RAG performs across
          the range of tasks an analyst works through, and which
          architectural choices drive that performance, remains
          empirically open.
    \item \textbf{G2 --- Evaluation realism:} Reported evaluations center
          on controlled benchmarks and isolated technical metrics, and
          ground-truth datasets for CTI tasks are scarce. Evidence that
          connects retrieval and generation quality to representative
          analyst tasks is rare.
    \item \textbf{G3 --- Source reliability:} RAG pipelines can ingest and
          propagate misinformation, and LLMs tend to over-trust retrieved
          context. How CTI-RAG systems should account for the integrity
          and reliability of their sources is not settled.
\end{enumerate}

These gaps, together with the workflow bottlenecks B1--B3 from
Section~\ref{sec:cti}, define the design space of this thesis:
Section~\ref{sec:requirements} derives design criteria from them, and
Section~\ref{sec:assessment} assesses the prototype against these
criteria.
